\renewcommand{\thetheorem}{12.\arabic{theorem}}
\setcounter{theorem}{0}
\setcounter{chapter}{11} % Sets chapter counter so the current chapter is 12

\chapter{Dimension and Subspaces}
In this chapter, we synthesize the fundamental results regarding finite-di\-men\-sional vector spaces and explore the relationship between a space and its subspaces through the lens of dimensionality.

\begin{summary}[Properties of Finite-Dimensional Spaces]
\label{sum:finite_dim_properties}
Let $V$ be a finite-di\-men\-sional vector space over a field $K$. The following structural principles hold:
\begin{enumerate}
    \item Every finite list of vectors that spans $V$ contains a sublist which constitutes a \newterm{basis} of $V$. Formally, if $(v_1, \dots, v_n)$ spans $V$, then there exists a sublist $(v_{i_1}, \dots, v_{i_d})$ with $1 \le i_1 < i_2 < \dots < i_d \le n$ that is a basis for $V$, where $d = \dim V$.
    \item Every linearly independent list of vectors in $V$ can be extended to a basis of $V$. If the list $(u_1, \dots, u_m)$ is linearly independent, there exist vectors $w_1, \dots, w_{d-m}$ such that 
    \[(u_1, \dots, u_m, w_1, \dots, w_{d-m})\] 
    is a basis of $V$, where $d = \dim V$.
\end{enumerate}
\end{summary}

% Creator: Opus 5
% Neither item is proved anywhere in the book, and item (2), the basis
% extension property, is used some ten times from ch. 14 onwards. Both follow
% in a few lines from the two lemmas of ch. 11. The derivation deliberately
% avoids \cref{thm:basis_equivalent_characterizations} below, whose own proof
% appeals to both items, which would make the argument circular.
The notes record these two principles as a summary of the previous lecture and
do not prove them. Everything in this chapter rests on them, and item
\textbf{(2)}, the \newterm{basis extension property}, is used repeatedly in the
chapters that follow, so we supply the derivation here. It is short, and it uses
nothing beyond the two lemmas of the previous chapter.

\begin{proof}[Proof of \cref{sum:finite_dim_properties}]
\textbf{(1)} This is \cref{lem:spanning_list_contains_basis}, which extracts a basis from within the list $(v_1, \dots, v_n)$ itself. Reading its elements in increasing order of index gives the sublist $(v_{i_1}, \dots, v_{i_d})$, and $d = \dim V$ by \cref{def:dimension}, since every basis of $V$ has that many elements.

\textbf{(2)} Let $(u_1, \dots, u_m)$ be linearly independent, and let $\mathcal{B} = (b_1, \dots, b_d)$ be a basis of $V$, which exists and has $d = \dim V$ elements by \cref{thm:basis_theorem}. In particular $\Sp(b_1, \dots, b_d) = V$, so \cref{lem:strong_steinitz} applies to the independent list $(u_1, \dots, u_m)$ and the spanning list $\mathcal{B}$. It yields $m \le d$ together with $d - m$ survivors $w_1, \dots, w_{d-m}$ of $\mathcal{B}$ such that
\[
\Sp(u_1, \dots, u_m, w_1, \dots, w_{d-m}) = V.
\]
This list spans $V$ and has exactly $d$ entries. By \cref{lem:spanning_list_contains_basis} it contains a sublist which is a basis of $V$, and that sublist has $d$ elements, again because every basis of $V$ does. A sublist of a list of length $d$ that itself has length $d$ is the whole list. Hence $(u_1, \dots, u_m, w_1, \dots, w_{d-m})$ is a basis of $V$, and it extends the given list.
\end{proof}
% Extractor: Gemini 3.1 Pro

% thm:12.1
\begin{theorem}[Equivalent Characterizations of a Basis]
\label{thm:basis_equivalent_characterizations}
Let $V$ be a vector space with $\dim V = n \in \mathbb{Z}_{\ge 0}$. Let $v_1, \dots, v_n \in V$ be a list of $n$ vectors. Then the following statements are equivalent:
\begin{enumerate}[label=\textbf{(\alph*)}]
    \item The vectors $v_1, \dots, v_n$ are linearly independent.
    \item The vectors $v_1, \dots, v_n$ generate $V$, i.e., $\Sp(v_1, \dots, v_n) = V$.
    \item The vectors $v_1, \dots, v_n$ constitute a basis for $V$.
\end{enumerate}
\end{theorem}

\begin{proof}
The equivalence is established through the following implications:

\textbf{(a) $\implies$ (c):} Suppose that \textbf{(a)} holds, so the list $v_1, \dots, v_n$ is linearly independent. According to the property that every linearly independent list can be extended to a basis, and observing that the list already possesses length $n = \dim V$, the extension requires no additional vectors. Thus, the list is already a basis for $V$, satisfying \textbf{(c)}.

\textbf{(c) $\implies$ (a) \& (b):} This follows directly from the definition of a basis, which by requirement must satisfy both the condition of linear independence in \textbf{(a)} and the spanning property in \textbf{(b)}.

\textbf{(b) $\implies$ (c):} If the spanning property in \textbf{(b)} holds, the list $v_1, \dots, v_n$ generates $V$. By the properties of finite-dimensional spaces, this list must contain a sublist that is a basis. Since $\dim V = n$, any basis must have exactly length $n$. Therefore, the sublist must be the entire list itself, making $v_1, \dots, v_n$ a basis as required by \textbf{(c)}.
\end{proof}



The previous theorem settles what happens for a list whose length is exactly $\dim V$. The next one records what happens when the length is wrong in either direction: too short a list cannot reach every vector, and too long a list must contain a redundancy.

\begin{theorem}[Lists of the Wrong Length]
\label{thm:wrong_length_lists}
Suppose that $n = \dim V < \infty$, and let $v_1, \dots, v_k$ be a list of vectors in $V$. Then:
\begin{enumerate}[label=\textbf{(\arabic*)}]
    \item If $k < n$, the list $v_1, \dots, v_k$ can \textbf{not} span $V$.
    \item If $k > n$, the list $v_1, \dots, v_k$ can \textbf{not} be linearly independent.
\end{enumerate}
\end{theorem}

\begin{proof}
\textbf{(1)} Assume that $k < n$, and suppose by contradiction that $\Sp(v_1, \dots, v_k) = V$. By item \textbf{(1)} of the summary on \cpageref{sum:finite_dim_properties} we can delete some of the vectors from the list $v_1, \dots, v_k$ and obtain a basis $v_{i_1}, \dots, v_{i_{k'}}$ of $V$. Then
\[
\dim V = k' \le k < n = \dim V,
\]
a contradiction. (The first equality is \cref{def:dimension}, applied to the basis just produced, and the second step holds because a sublist is no longer than the list it comes from.)

\textbf{(2)} Assume that $k > n$, and suppose by contradiction that $v_1, \dots, v_k$ are linearly independent. By item \textbf{(2)} of the summary on \cpageref{sum:finite_dim_properties} we can extend this list to a basis of $V$. Then
\[
\dim V \ge k > n = \dim V,
\]
a contradiction. (The first inequality holds because that basis contains all $k$ vectors of the list, and has $\dim V$ elements by \cref{def:dimension}.)
\end{proof}

% prop:12.2
\begin{proposition}[Dimensional Constraints on Subspaces]
\label{prop:subspace_dimension_constraints}
Let $V$ be a finite-di\-men\-sional vector space over $K$. Then every linear subspace $U \subseteq V$ is also finite-di\-men\-sional. Moreover, $\dim U \le \dim V$, with the equality $\dim U = \dim V$ holding if and only if $U = V$.
\end{proposition}

\begin{proof}
Let $n := \dim V$. We first establish that $U$ possesses a finite basis. If $U = \{0\}$, the empty set serves as the basis and we are done. Assume $U \neq \{0\}$. We select a non-zero vector $v_1 \in U$. If $\Sp(v_1) = U$, we have found our basis. 

If $\Sp(v_1) \neq U$, we choose a second vector $v_2 \in U \setminus \Sp(v_1)$. It follows from the properties of linear independence (\cref{lem:no_predecessor_span}) that the list $\{v_1, v_2\}$ is linearly independent. We proceed iteratively: after $j$ steps, we obtain a linearly independent list $\{v_1, \dots, v_j\} \subseteq U$. If $\Sp(v_1, \dots, v_j) = U$, the process terminates. Otherwise, we choose $v_{j+1} \in U \setminus \Sp(v_1, \dots, v_j)$. By the principle that adding a vector outside the span of a linearly independent set preserves independence (again \cref{lem:no_predecessor_span}), the expanded list $\{v_1, \dots, v_{j+1}\}$ remains linearly independent.

Crucially, in a space $V$ of dimension $n$, no linearly independent list can exceed length $n$; this is \cref{thm:wrong_length_lists} \textbf{(2)}. Consequently, this algorithmic procedure must terminate in $k$ steps, where $k \le n$. At this stage, we have a list $\{v_1, \dots, v_k\}$ that is linearly independent and spans $U$, thereby constituting a basis for $U$. This implies $U$ is finite-dimensional and $\dim U = k \le n = \dim V$.

It remains to settle when equality holds. One direction is immediate: if $U = V$, then clearly $\dim U = \dim V$. The converse is the substantial one.

Conversely, if $\dim U = \dim V = n$, then $U$ contains a linearly independent list of length $n$, namely a basis $v_1, \dots, v_n$ of $U$, which stays independent when read as a list in $V$, independence being a condition on linear combinations and not on the ambient space. By \Cref{thm:basis_equivalent_characterizations}, such a list satisfies property \textbf{(a)} and is therefore also a basis for $V$. This implies $\Sp(v_1, \dots, v_n) = V$. Since $U = \Sp(v_1, \dots, v_n)$ by the construction of the basis, we conclude $U = V$.
\end{proof}

\begin{exercise}[Dimension and Bases of Matrix Subspaces]
\label{exc:matrix_subspaces_dimension_bases}
Compute the dimension and find an explicit basis for each of the following vector spaces over the given field:
\begin{enumerate}
    \item The space of upper triangular matrices in $\M_{n \times n}(\mathbb{R})$ over $\mathbb{R}$.
    \item The space of diagonal matrices in $\M_{n \times n}(\mathbb{R})$ over $\mathbb{R}$.
    \item The space of symmetric matrices
    \[
    W := \{A \in \M_{n \times n}(\mathbb{R}) \mid A\transp = A\} \subseteq \M_{n \times n}(\mathbb{R}).
    \]
    \item The space of matrices $A \in \M_{n \times n}(\mathbb{F}_2)$ such that the sum of the columns of $A$ is the zero vector in $\mathbb{F}_2^n$.
\end{enumerate}
\exinfo{This exercise is Problem 1 of Problem Sheet 8, Linear Algebra I, Autumn Semester 2025.}
\exsol{sol:matrix_subspaces_dimension_bases}
\end{exercise}

\begin{exercise}[Subspace Dimension, Defining Equations, and Basis Selection]
\label{exc:subspace_dimension_equations_basis_selection}
Let $W \subseteq \mathbb{R}^4$ be the linear subspace spanned by the four vectors
\[
v_1 = \begin{pmatrix} 1 \\ 2 \\ 1 \\ 2 \end{pmatrix}, \qquad
v_2 = \begin{pmatrix} 0 \\ 0 \\ -1 \\ 1 \end{pmatrix}, \qquad
v_3 = \begin{pmatrix} 1 \\ 4 \\ 2 \\ 3 \end{pmatrix}, \qquad
v_4 = \begin{pmatrix} 0 \\ 2 \\ 1 \\ 1 \end{pmatrix}.
\]
\begin{enumerate}
    \item Compute the dimension $\dim W$.
    \item Determine a homogeneous system of linear equations whose set of solutions is precisely $W$.
    \item Find all possible subsets of $\{v_1, v_2, v_3, v_4\}$ that form a basis of $W$. How many such bases are there?
\end{enumerate}
\exhint[Official Hint]{For part (b), if $\dim W = 3$, find a non-zero normal vector $n = (n_1, n_2, n_3, n_4)\transp \in \mathbb{R}^4$ orthogonal to all $w \in W$ (i.e., $\sum_{i=1}^4 n_i w_i = 0$ for all $w \in W$). Then $W$ is described by the single equation $n_1 x_1 + n_2 x_2 + n_3 x_3 + n_4 x_4 = 0$.}
\exinfo{This exercise is Problem 2 of Problem Sheet 8, Linear Algebra I, Autumn Semester 2025.}
\exsol{sol:subspace_dimension_equations_basis_selection}
\end{exercise}

\begin{exercise}[Bases and Dimensions of Various Spaces]
\label{exc:bases_dimensions_various_spaces}
Determine an explicit basis and the dimension of each of the following vector spaces:
\begin{enumerate}
    \item The solution set $S \subseteq \mathbb{R}^3$ of the homogeneous linear system:
    \[
    \begin{aligned}
    x + y - z &= 0, \\
    3x + y + 2z &= 0, \\
    2x + 3z &= 0.
    \end{aligned}
    \]
    \item The trivial vector space $\{0\}$.
    \item The space $V := \{(z, w) \in \mathbb{C}^2 \mid z + iw = 0\}$ viewed as a vector space over $\mathbb{C}$.
    \item The same space $V = \{(z, w) \in \mathbb{C}^2 \mid z + iw = 0\}$ viewed as a vector space over $\mathbb{R}$.
\end{enumerate}
\exinfo{This exercise is Problem 3 of Problem Sheet 8, Linear Algebra I, Autumn Semester 2025.}
\exsol{sol:bases_dimensions_various_spaces}
\end{exercise}

\begin{exercise}[Dimension of Polynomials with a Linear Factor]
\label{exc:polynomials_linear_factor_dimension}
Let $K$ be a field, fix the polynomial $g(X) := X + 5 \in K[X]$, and let $d \ge 1$. Compute the dimension of the subspace
\[
W := \{h \in K[X] \mid \deg(h) \le d \text{ and } \exists f \in K[X] \text{ such that } h = g \cdot f\} \subseteq K[X]_d.
\]
\exinfo{This exercise is Problem 4 of Problem Sheet 8, Linear Algebra I, Autumn Semester 2025.}
\exsol{sol:polynomials_linear_factor_dimension}
\end{exercise}

\begin{exercise}[Basis and Dimension of Second-Order Harmonic ODE]
\label{exc:harmonic_ode_solution_space}
Determine an explicit basis and compute the dimension of the real vector space
\[
V := \left\{ f \colon \mathbb{R} \to \mathbb{R} \;\middle|\; f \text{ is twice differentiable and } f'' + f = 0 \right\}.
\]
\exhint[Official Hint]{Recall from analysis that the only continuous solution to the initial value problem $y'' + y = 0$ with $y(0) = 0$ and $y'(0) = 0$ is the identically zero function $y \equiv 0$. Consider evaluating functions and their first derivatives at $0$.}
\exinfo{This exercise is Problem 5 of Problem Sheet 8, Linear Algebra I, Autumn Semester 2025.}
\exsol{sol:harmonic_ode_solution_space}
\end{exercise}

\begin{exercise}[Dimension of a Shifted Span]
\label{exc:shifted_span_dimension}
Suppose that $\{v_1, \dots, v_m\}$ is a linearly independent subset of a vector space $V$, and let $w \in V$ be an arbitrary vector. Prove that
\[
\dim \Sp(v_1 + w, v_2 + w, \dots, v_m + w) \ge m - 1.
\]
\exinfo{This exercise is Problem 4 of Problem Sheet 9, Linear Algebra I, Autumn Semester 2025.}
\exsol{sol:shifted_span_dimension}
\end{exercise}

\begin{exercise}[Descending Chains of Subspaces]
\label{exc:descending_chains_subspaces}
Let $V$ be a vector space over a field $F$, and let
\[
V \supseteq U_0 \supseteq U_1 \supseteq U_2 \supseteq \dots \supseteq U_k \supseteq \dots
\]
be a sequence of nested linear subspaces.
\begin{enumerate}
    \item If $V$ is finite-dimensional, prove that the sequence stabilises: there exists $N \in \mathbb{N}$ such that $U_n = U_N$ for all $n \ge N$.
    \item Is this necessarily true if $V$ is infinite-dimensional?
    \item Assume that $V$ is infinite-dimensional and that $\dim U_n \ge 1$ for all $n \in \mathbb{N}$. What can you say about $\bigcap_{n=0}^\infty U_n$? Can this intersection be $\{0\}$?
\end{enumerate}
\exinfo{This exercise is Problem 6 of Problem Sheet 9, Linear Algebra I, Autumn Semester 2025.}
\exsol{sol:descending_chains_subspaces}
\end{exercise}

\begin{ainote}
\Cref{thm:wrong_length_lists}, and the trivial direction of the equality claim in
\cref{prop:subspace_dimension_constraints}, were missing from this transcript and
are restored above. The notes also restate \cref{lem:no_predecessor_span} in full
at this point, as a recollection from the earlier lecture; here it is only
cross-referenced.

The substantial addition is the proof of \cref{sum:finite_dim_properties}. The
notes give its two principles as a recollection and prove neither, and nothing
else in the book proves them, although item \textbf{(2)}, the basis extension
property, is used some ten times from ch. 14 onwards and is occasionally credited
in passing to \qt{the Steinitz Exchange Lemma}. Both follow from
\cref{lem:spanning_list_contains_basis} and \cref{lem:strong_steinitz} in a few
lines. The derivation deliberately avoids
\cref{thm:basis_equivalent_characterizations}, which stands just below it and
whose own proof appeals to both items: going through it would be circular. That
is the reason for the longer route, and a later simplification should not undo
it.
\end{ainote}

\section{Solutions}
\label{sec:solutions_dimension}

% Official Solution (Problem Sheet 8, Exercise 1)
\begin{exercisesolution}[Solution to \cref{exc:matrix_subspaces_dimension_bases}]
\label{sol:matrix_subspaces_dimension_bases}
Let $E_{ij} \in \M_{n \times n}(\mathbb{R})$ denote the standard elementary matrix having entry $1$ in the $(i, j)$-position and $0$ elsewhere, for $1 \le i, j \le n$. The set of all $n^2$ elementary matrices forms the standard basis of $\M_{n \times n}(\mathbb{R})$.

\begin{enumerate}[label=\textbf{(\alph*)}]
    \item \textbf{Upper triangular matrices.} A matrix $A = (a_{ij}) \in \M_{n \times n}(\mathbb{R})$ is upper triangular if and only if $a_{ij} = 0$ whenever $i > j$. Any such matrix can be expressed uniquely as
    \[
        A = \sum_{1 \le i \le j \le n} a_{ij} E_{ij}.
    \]
    The family $\mathcal{B} := \{E_{ij} \mid 1 \le i \le j \le n\}$ is a subset of the standard basis, hence linearly independent, and it spans the subspace of upper triangular matrices. Thus, $\mathcal{B}$ is a basis. The number of basis elements is
    \[
        \sum_{i=1}^n (n - i + 1) = \sum_{k=1}^n k = \frac{n(n+1)}{2}.
    \]
    Consequently, the dimension of the space is $\frac{n(n+1)}{2}$.

    \item \textbf{Diagonal matrices.} A matrix $A = (a_{ij})$ is diagonal if and only if $a_{ij} = 0$ for all $i \neq j$. Every diagonal matrix can be written uniquely as
    \[
        A = \sum_{i=1}^n a_{ii} E_{ii}.
    \]
    The set $\mathcal{B} := \{E_{ii} \mid 1 \le i \le n\}$ is linearly independent and spans the subspace of diagonal matrices. Hence, $\mathcal{B}$ is a basis, and the dimension is $n$.

    \item \textbf{Symmetric matrices.} A matrix $A = (a_{ij})$ is symmetric if and only if $a_{ji} = a_{ij}$ for all $1 \le i, j \le n$. For $1 \le i \le n$, set $D_i := E_{ii}$, and for $1 \le i < j \le n$, set $S_{ij} := E_{ij} + E_{ji}$. Each $D_i$ and $S_{ij}$ is symmetric since $D_i\transp = D_i$ and $S_{ij}\transp = (E_{ij} + E_{ji})\transp = E_{ji} + E_{ij} = S_{ij}$.
    
    Any symmetric matrix $A \in W$ satisfies
    \[
        A = \sum_{i=1}^n a_{ii} D_i + \sum_{1 \le i < j \le n} a_{ij} S_{ij}.
    \]
    To verify linear independence, suppose that
    \[
        \sum_{i=1}^n \alpha_{ii} D_i + \sum_{1 \le i < j \le n} \alpha_{ij} S_{ij} = 0.
    \]
    Examining the diagonal entries yields $\alpha_{ii} = 0$ for each $1 \le i \le n$. For $i < j$, the $(i, j)$-entry of the linear combination is $\alpha_{ij}$, which implies $\alpha_{ij} = 0$. Thus, the family
    \[
        \mathcal{B} := \{D_i \mid 1 \le i \le n\} \cup \{S_{ij} \mid 1 \le i < j \le n\}
    \]
    is linearly independent and spans $W$, so it is a basis. Its cardinality is
    \[
        n + \binom{n}{2} = n + \frac{n(n-1)}{2} = \frac{n(n+1)}{2}.
    \]
    Therefore, $\dim W = \frac{n(n+1)}{2}$.

    \item \textbf{Matrices over $\mathbb{F}_2$ whose columns sum to the zero vector.} Let $W \subseteq \M_{n \times n}(\mathbb{F}_2)$ denote this subspace. The sum of the columns of $A$ is the vector whose $i$-th entry is $\sum_{j=1}^n a_{ij}$, so the condition is read off row by row: a matrix $A = (a_{ij}) \in \M_{n \times n}(\mathbb{F}_2)$ belongs to $W$ if and only if for each row index $1 \le i \le n$, the sum of the entries across that row vanishes in $\mathbb{F}_2$:
    \[
        \sum_{j=1}^n a_{ij} = 0 \iff a_{in} = \sum_{j=1}^{n-1} a_{ij} \quad \text{in } \mathbb{F}_2.
    \]
    Thus, for each of the $n$ rows, the first $n-1$ entries can be chosen arbitrarily, while the $n$-th entry is uniquely determined by the zero-sum condition.
    
    For each $1 \le i \le n$ and $1 \le k \le n-1$, define the matrix
    \[
        A_k^{(i)} := E_{ik} + E_{in} \in \M_{n \times n}(\mathbb{F}_2).
    \]
    In $A_k^{(i)}$, row $i$ has $1$ in column $k$ and $1$ in column $n$ (whose sum is $1 + 1 = 0 \in \mathbb{F}_2$), and all other rows are zero, so $A_k^{(i)} \in W$. Any $A \in W$ can be decomposed as
    \[
        A = \sum_{i=1}^n \sum_{k=1}^{n-1} a_{ik} A_k^{(i)}.
    \]
    To verify linear independence, suppose that $\sum_{i=1}^n \sum_{k=1}^{n-1} \lambda_{ik} A_k^{(i)} = 0$. For any fixed pair $(i, k)$ with $1 \le k \le n-1$, the $(i, k)$-entry of the sum is precisely $\lambda_{ik}$, forcing $\lambda_{ik} = 0$.
    
    Thus the collection $\mathcal{B} := \{A_k^{(i)} \mid 1 \le i \le n,\, 1 \le k \le n-1\}$ is a basis of $W$, and its cardinality is $n(n-1)$. Hence, $\dim W = n(n-1)$.
\end{enumerate}
\exinfo{Official Solution (Problem Sheet 8, Exercise 1). Determines bases and dimensions of upper triangular, diagonal, symmetric, and row-sum zero matrix spaces.}
\end{exercisesolution}

% Official Solution (Problem Sheet 8, Exercise 2)
\begin{exercisesolution}[Solution to \cref{exc:subspace_dimension_equations_basis_selection}]
\label{sol:subspace_dimension_equations_basis_selection}
\begin{enumerate}[label=\textbf{(\alph*)}]
    \item We determine the dimension of $W = \Sp(v_1, v_2, v_3, v_4)$ by applying elementary row operations to the matrix whose columns are the given vectors:
    \[
        M := \begin{pmatrix}
            1 & 0 & 1 & 0 \\
            2 & 0 & 4 & 2 \\
            1 & -1 & 2 & 1 \\
            2 & 1 & 3 & 1
        \end{pmatrix}.
    \]
    Subtracting multiples of the first row from the lower rows via $R_2 - 2R_1 \to R_2$, $R_3 - R_1 \to R_3$, and $R_4 - 2R_1 \to R_4$ yields
    \[
        \begin{pmatrix}
            1 & 0 & 1 & 0 \\
            0 & 0 & 2 & 2 \\
            0 & -1 & 1 & 1 \\
            0 & 1 & 1 & 1
        \end{pmatrix}.
    \]
    Swapping $R_2 \leftrightarrow R_4$ produces
    \[
        \begin{pmatrix}
            1 & 0 & 1 & 0 \\
            0 & 1 & 1 & 1 \\
            0 & -1 & 1 & 1 \\
            0 & 0 & 2 & 2
        \end{pmatrix}.
    \]
    Adding the second row to the third row ($R_3 + R_2 \to R_3$) gives
    \[
        \begin{pmatrix}
            1 & 0 & 1 & 0 \\
            0 & 1 & 1 & 1 \\
            0 & 0 & 2 & 2 \\
            0 & 0 & 2 & 2
        \end{pmatrix},
    \]
    and subtracting the third row from the fourth row ($R_4 - R_3 \to R_4$) results in the row echelon form
    \[
        \begin{pmatrix}
            1 & 0 & 1 & 0 \\
            0 & 1 & 1 & 1 \\
            0 & 0 & 2 & 2 \\
            0 & 0 & 0 & 0
        \end{pmatrix}.
    \]
    There are $3$ non-zero pivot rows, so the column rank is $3$. Therefore, $\dim W = 3$.

    \item Since $\dim W = 3$ inside $\mathbb{R}^4$, the subspace $W$ is a hyperplane and can be defined by a single homogeneous linear equation $n_1 x_1 + n_2 x_2 + n_3 x_3 + n_4 x_4 = 0$, where $n = (n_1, n_2, n_3, n_4)\transp \in \mathbb{R}^4 \setminus \{0\}$ is orthogonal to every vector in $W$.
    
    It suffices to enforce orthogonality on the spanning vectors $v_1, v_2, v_4$:
    \[
        \begin{aligned}
            n_1 + 2n_2 + n_3 + 2n_4 &= 0, \\
            -n_3 + n_4 &= 0, \\
            2n_2 + n_3 + n_4 &= 0.
        \end{aligned}
    \]
    The second equation gives $n_4 = n_3$. Substituting this into the third equation yields $2n_2 + 2n_3 = 0$, so $n_2 = -n_3$. Substituting both into the first equation gives
    \[
        n_1 + 2(-n_3) + n_3 + 2n_3 = n_1 + n_3 = 0 \implies n_1 = -n_3.
    \]
    Setting $n_3 := -1$ gives $n_1 = 1$, $n_2 = 1$, and $n_4 = -1$, yielding the normal vector $n = (1, 1, -1, -1)\transp$. We readily check that $\langle n, v_3 \rangle = 1(1) + 1(4) - 1(2) - 1(3) = 0$.
    
    Therefore, the subspace $W$ is precisely the solution set of the single equation
    \[
        x_1 + x_2 - x_3 - x_4 = 0.
    \]

    \item Since $\dim W = 3$, any basis chosen from $\{v_1, v_2, v_3, v_4\}$ must consist of exactly $3$ vectors. There are $\binom{4}{3} = 4$ three-element subsets:
    \begin{itemize}
        \item Observe that $v_3 = v_1 + v_4$, since
        \[
            \begin{pmatrix} 1 \\ 2 \\ 1 \\ 2 \end{pmatrix} + \begin{pmatrix} 0 \\ 2 \\ 1 \\ 1 \end{pmatrix} = \begin{pmatrix} 1 \\ 4 \\ 2 \\ 3 \end{pmatrix}.
        \]
        Thus the subset $\{v_1, v_3, v_4\}$ is linearly dependent and spans only a $2$-dimensional subspace; it is not a basis.
        \item For $\{v_1, v_2, v_4\}$, the Gaussian elimination in part \textbf{(a)} established that the corresponding columns are linearly independent. By \cref{thm:basis_equivalent_characterizations}, any linearly independent list of length $3$ in a $3$-dimensional space is a basis.
        \item For $\{v_1, v_2, v_3\}$, since $v_4 = v_3 - v_1$, we have $\Sp(v_1, v_2, v_3) = \Sp(v_1, v_2, v_3, v_4) = W$. By \cref{thm:basis_equivalent_characterizations}, this spanning set of size $3$ is a basis.
        \item For $\{v_2, v_3, v_4\}$, since $v_1 = v_3 - v_4$, we have $\Sp(v_2, v_3, v_4) = W$, so this is likewise a basis.
    \end{itemize}
    Thus there are exactly $3$ such bases:
    \[
        \mathcal{B}_1 := \{v_1, v_2, v_3\}, \qquad \mathcal{B}_2 := \{v_1, v_2, v_4\}, \qquad \mathcal{B}_3 := \{v_2, v_3, v_4\}.
    \]
\end{enumerate}
\exinfo{Official Solution (Problem Sheet 8, Exercise 2). Computes basis selection, dimension, and defining homogeneous Cartesian equations for a subspace in $\mathbb{R}^4$.}
\end{exercisesolution}

% Solution: Gemini / Official Hybrid (Problem Sheet 8, Exercise 3)
\begin{exercisesolution}[Solution to \cref{exc:bases_dimensions_various_spaces}]
\label{sol:bases_dimensions_various_spaces}
\begin{enumerate}[label=\textbf{(\alph*)}]
    \item \textbf{System of linear equations.} The system is
    \[
        \begin{aligned}
            x + y - z &= 0, \\
            3x + y + 2z &= 0, \\
            2x + 3z &= 0.
        \end{aligned}
    \]
    Row reducing the coefficient matrix:
    \[
        \begin{pmatrix}
            1 & 1 & -1 \\
            3 & 1 & 2 \\
            2 & 0 & 3
        \end{pmatrix}
        \xrightarrow{\substack{R_2 - 3R_1 \to R_2 \\ R_3 - 2R_1 \to R_3}}
        \begin{pmatrix}
            1 & 1 & -1 \\
            0 & -2 & 5 \\
            0 & -2 & 5
        \end{pmatrix}
        \xrightarrow{R_3 - R_2 \to R_3}
        \begin{pmatrix}
            1 & 1 & -1 \\
            0 & -2 & 5 \\
            0 & 0 & 0
        \end{pmatrix}.
    \]
    The second row gives $-2y + 5z = 0 \implies y = \frac{5}{2}z$. Substituting into the first row yields $x = z - y = z - \frac{5}{2}z = -\frac{3}{2}z$. Choosing $z := 2$ gives the solution vector $(-3, 5, 2)\transp$.
    
    Thus the solution set is $S = \Sp\left( \begin{pmatrix} -3 \\ 5 \\ 2 \end{pmatrix} \right)$. Since the vector is non-zero, $\mathcal{B} := \left\{ \begin{pmatrix} -3 \\ 5 \\ 2 \end{pmatrix} \right\}$ is a basis of $S$, and $\dim S = 1$.

    \item \textbf{Trivial space $\{0\}$.} The empty set $\emptyset$ is linearly independent and spans $\{0\}$ (because the empty sum of vectors evaluates to $0$). Thus, $\mathcal{B} := \emptyset$ is a basis of $\{0\}$, and $\dim(\{0\}) = |\emptyset| = 0$.

    \item \textbf{$V = \{(z, w) \in \mathbb{C}^2 \mid z + iw = 0\}$ over $\mathbb{C}$.} The condition $z + iw = 0$ is equivalent to $z = -iw$. Every vector in $V$ is of the form
    \[
        \begin{pmatrix} z \\ w \end{pmatrix} = \begin{pmatrix} -iw \\ w \end{pmatrix} = w \begin{pmatrix} -i \\ 1 \end{pmatrix} \quad (w \in \mathbb{C}).
    \]
    The single vector $v_1 := (-i, 1)\transp \in \mathbb{C}^2$ is non-zero, hence linearly independent over $\mathbb{C}$, and it generates $V$. Thus, $\mathcal{B} := \left\{ \begin{pmatrix} -i \\ 1 \end{pmatrix} \right\}$ is a basis of $V$ over $\mathbb{C}$, and $\dim_{\mathbb{C}} V = 1$.

    \item \textbf{$V = \{(z, w) \in \mathbb{C}^2 \mid z + iw = 0\}$ over $\mathbb{R}$.} Writing the free parameter $w \in \mathbb{C}$ in terms of real and imaginary parts as $w = a + ib$ with $a, b \in \mathbb{R}$, we obtain
    \[
        \begin{pmatrix} z \\ w \end{pmatrix} = (a + ib) \begin{pmatrix} -i \\ 1 \end{pmatrix} = a \begin{pmatrix} -i \\ 1 \end{pmatrix} + b \begin{pmatrix} 1 \\ i \end{pmatrix}.
    \]
    Let $u_1 := (-i, 1)\transp$ and $u_2 := (1, i)\transp$. These two vectors belong to $V$ and generate $V$ as a real vector space. To verify linear independence over $\mathbb{R}$, suppose that $a u_1 + b u_2 = 0$ for $a, b \in \mathbb{R}$. Looking at the second coordinate gives
    \[
        a(1) + b(i) = a + ib = 0 \implies a = 0 \quad \text{and} \quad b = 0.
    \]
    Hence $\mathcal{B} := \left\{ \begin{pmatrix} -i \\ 1 \end{pmatrix}, \begin{pmatrix} 1 \\ i \end{pmatrix} \right\}$ is a basis of $V$ over $\mathbb{R}$, and $\dim_{\mathbb{R}} V = 2$.
\end{enumerate}
\exinfo{Hybrid Solution (Gemini 3.7 Flash / Official Problem Sheet 8, Exercise 3). Finds bases and dimensions for linear systems, trivial space, and complex vs real planes.}
\end{exercisesolution}

% Solution: Gemini / Official Hybrid (Problem Sheet 8, Exercise 4)
\begin{exercisesolution}[Solution to \cref{exc:polynomials_linear_factor_dimension}]
\label{sol:polynomials_linear_factor_dimension}
Let $g(X) := X + 5 \in K[X]$. An element $h \in K[X]$ belongs to $W$ if and only if $\deg(h) \le d$ and $h(X) = g(X) f(X)$ for some polynomial $f(X) \in K[X]$.

For any non-zero polynomial $f \in K[X]$, the degree formula yields
\[
    \deg(h) = \deg(g \cdot f) = \deg(g) + \deg(f) = 1 + \deg(f).
\]
Therefore, $\deg(h) \le d$ is equivalent to $\deg(f) \le d - 1$. The subspace of polynomials of degree at most $d - 1$ is spanned by the standard monomial basis $\{1, X, X^2, \dots, X^{d-1}\}$. Writing $f(X) = \sum_{j=0}^{d-1} a_j X^j$ with coefficients $a_j \in K$, we obtain
\[
    h(X) = g(X) \sum_{j=0}^{d-1} a_j X^j = \sum_{j=0}^{d-1} a_j X^j g(X).
\]
This shows that the set of $d$ polynomials
\[
    \mathcal{B} := \{g(X), X g(X), X^2 g(X), \dots, X^{d-1} g(X)\}
\]
spans $W$.

To prove linear independence, consider a vanishing linear combination
\[
    \sum_{j=0}^{d-1} c_j X^j g(X) = 0 \iff g(X) \cdot \left( \sum_{j=0}^{d-1} c_j X^j \right) = 0.
\]
Since $K[X]$ is an integral domain and $g(X) = X + 5 \neq 0$, the second factor must vanish identically:
\[
    \sum_{j=0}^{d-1} c_j X^j = 0.
\]
Because the standard monomials $1, X, \dots, X^{d-1}$ are linearly independent in $K[X]$, it follows that $c_0 = c_1 = \dots = c_{d-1} = 0$.

Thus, $\mathcal{B}$ is a basis of $W$. Since it contains exactly $d$ polynomials, we conclude that $\dim W = d$.
\exinfo{Hybrid Solution (Gemini 3.7 Flash / Official Problem Sheet 8, Exercise 4). Computes the dimension of multiples of a fixed linear factor in polynomial spaces.}
\end{exercisesolution}

% Solution: Gemini / Official Hybrid (Problem Sheet 8, Exercise 5)
\begin{exercisesolution}[Solution to \cref{exc:harmonic_ode_solution_space}]
\label{sol:harmonic_ode_solution_space}
We determine the dimension and an explicit basis for the real vector space
\[
    V := \left\{ f \colon \mathbb{R} \to \mathbb{R} \;\middle|\; f \text{ is twice differentiable and } f'' + f = 0 \right\}.
\]

\textbf{Step 1: Two particular solutions.} The functions $s(x) := \sin x$ and $c(x) := \cos x$ are infinitely differentiable and satisfy $s''(x) = -\sin x = -s(x)$ and $c''(x) = -\cos x = -c(x)$, so $s, c \in V$.

\textbf{Step 2: The spanning property.} Let $f \in V$ be arbitrary. Set $a := f(0) \in \mathbb{R}$ and $b := f'(0) \in \mathbb{R}$. Define the auxiliary function
\[
    g(x) := f(x) - a \cos x - b \sin x.
\]
Because the differential operator $L(f) := f'' + f$ is linear, the difference of solutions is again a solution, so $g \in V$, meaning $g'' + g = 0$. Evaluating $g$ and its first derivative at $x = 0$:
\[
    g(0) = f(0) - a \cos 0 - b \sin 0 = a - a(1) - b(0) = 0,
\]
\[
    g'(0) = f'(0) + a \sin 0 - b \cos 0 = b + a(0) - b(1) = 0.
\]
According to the uniqueness theorem for linear second-order ordinary differential equations (as recalled in the hint), the only solution to the initial value problem $y'' + y = 0$ with $y(0) = 0$ and $y'(0) = 0$ is the identically zero function $y \equiv 0$.
Consequently, $g(x) = 0$ for all $x \in \mathbb{R}$, which yields
\[
    f(x) = a \cos x + b \sin x.
\]
This proves that $\{\cos x, \sin x\}$ spans $V$.

\textbf{Step 3: Linear independence.} Suppose that $\lambda_1 \cos x + \lambda_2 \sin x = 0$ for all $x \in \mathbb{R}$, where $\lambda_1, \lambda_2 \in \mathbb{R}$. Evaluating at $x = 0$ gives $\lambda_1(1) + \lambda_2(0) = \lambda_1 = 0$. Evaluating at $x = \pi/2$ gives $\lambda_1(0) + \lambda_2(1) = \lambda_2 = 0$. Thus, $\cos x$ and $\sin x$ are linearly independent over $\mathbb{R}$.

\textbf{Conclusion.} The set $\mathcal{B} := \{\cos x, \sin x\}$ is an explicit basis of $V$, and $\dim V = 2$.
\exinfo{Hybrid Solution (Gemini 3.7 Flash / Official Problem Sheet 8, Exercise 5). Demonstrates the 2-dimensional solution space of harmonic oscillator ODE $f'' + f = 0$.}
\end{exercisesolution}

% Solution: Gemini / Official Hybrid (Problem Sheet 9, Exercise 4)
\begin{exercisesolution}[Proof of \cref{exc:shifted_span_dimension}]
\label{sol:shifted_span_dimension}
Let $\{v_1, \dots, v_m\}$ be a linearly independent subset of $V$, let $w \in V$, and define the subspace
\[
    W := \Sp(v_1 + w, v_2 + w, \dots, v_m + w) \subseteq V.
\]
For each index $i \in \{2, \dots, m\}$, consider the difference of generators:
\[
    u_i := (v_i + w) - (v_1 + w) = v_i - v_1.
\]
Since both $v_i + w$ and $v_1 + w$ belong to $W$, their difference $u_i = v_i - v_1$ belongs to $W$ for all $2 \le i \le m$.

We show that the list of $m - 1$ vectors $\{u_2, \dots, u_m\} = \{v_2 - v_1, \dots, v_m - v_1\}$ is linearly independent. Suppose that $\alpha_2, \dots, \alpha_m \in K$ are scalars such that
\[
    \sum_{i=2}^m \alpha_i u_i = 0 \iff \sum_{i=2}^m \alpha_i (v_i - v_1) = 0.
\]
Rearranging terms by grouping the vectors $v_1, \dots, v_m$ gives
\[
    \left( -\sum_{i=2}^m \alpha_i \right) v_1 + \sum_{i=2}^m \alpha_i v_i = 0.
\]
Because the original family $\{v_1, v_2, \dots, v_m\}$ is linearly independent, every scalar in this linear combination must vanish. In particular, examining the coefficient of $v_i$ for each $2 \le i \le m$ yields
\[
    \alpha_i = 0 \quad \text{for all } 2 \le i \le m.
\]
Hence, $\{u_2, \dots, u_m\} \subseteq W$ is a linearly independent subset of cardinality $m - 1$. By \cref{prop:subspace_dimension_constraints}, every subspace containing a linearly independent set of $m - 1$ elements must have dimension at least $m - 1$. Therefore,
\[
    \dim \Sp(v_1 + w, v_2 + w, \dots, v_m + w) \ge m - 1.
\]
\exinfo{Hybrid Solution (Gemini 3.7 Flash / Official Problem Sheet 9, Exercise 4). Proves the dimension bound $\dim \Sp(v_1+w, \dots, v_m+w) \ge m-1$ via difference vectors.}
\end{exercisesolution}

% Official Solution (Problem Sheet 9, Exercise 6)
\begin{exercisesolution}[Solution to \cref{exc:descending_chains_subspaces}]
\label{sol:descending_chains_subspaces}
\begin{enumerate}[label=\textbf{(\alph*)}]
    \item \textbf{Finite-dimensional case.} Suppose that $\dim V = d < \infty$. By \cref{prop:subspace_dimension_constraints}, every linear subspace of $V$ is finite-dimensional with $\dim U \le \dim V$. Since the sequence is nested with $U_{n+1} \subseteq U_n$, applying \cref{prop:subspace_dimension_constraints} gives $\dim U_{n+1} \le \dim U_n$ for all $n \ge 0$.
    
    Thus the sequence of dimensions
    \[
        d \ge \dim U_0 \ge \dim U_1 \ge \dim U_2 \ge \dots \ge 0
    \]
    is a non-increasing sequence of non-negative integers. Because a non-increasing sequence in $\mathbb{Z}_{\ge 0}$ cannot decrease infinitely often, it must eventually become constant: there exist $N \in \mathbb{N}$ and an integer $k \in \{0, \dots, d\}$ such that
    \[
        \dim U_n = k \quad \text{for all } n \ge N.
    \]
    For any $n \ge N$, we have $U_{n+1} \subseteq U_n$ and $\dim U_{n+1} = \dim U_n = k$. By the equality criterion of \cref{prop:subspace_dimension_constraints}, an inclusion of finite-dimensional subspaces with equal dimensions forces equality of the subspaces:
    \[
        U_{n+1} = U_n \quad \text{for all } n \ge N.
    \]
    By induction, $U_n = U_N$ for all $n \ge N$, so the sequence stabilizes.

    \item \textbf{Infinite-dimensional case.} This is not necessarily true if $V$ is infinite-dimensional. Consider the polynomial space $V := F[X]$ over the field $F$. For each $n \ge 0$, define the subspace
    \[
        U_n := \Sp(X^n, X^{n+1}, X^{n+2}, \dots) = \{p(X) \in F[X] \mid X^n \text{ divides } p(X)\}.
    \]
    Clearly $U_{n+1} \subseteq U_n$ for all $n \ge 0$. However, the monomial $X^n$ belongs to $U_n$ but $X^n \notin U_{n+1}$, so $U_{n+1} \subsetneq U_n$ is a strictly proper inclusion for every $n \ge 0$. Consequently, the sequence never stabilizes.

    \item \textbf{Infinite intersection.} If $V$ is infinite-dimensional and $\dim U_n \ge 1$ for all $n \in \mathbb{N}$, the intersection $\bigcap_{n=0}^\infty U_n$ can be $\{0\}$.
    
    To see this, consider again the nested sequence $U_n = \Sp(X^n, X^{n+1}, \dots) \subset F[X]$ from part \textbf{(b)}. Each subspace $U_n$ contains the infinite linearly independent set $\{X^k \mid k \ge n\}$, so $\dim U_n = \infty \ge 1$ for all $n$.
    
    Now suppose $p(X) \in \bigcap_{n=0}^\infty U_n$. If $p(X) \neq 0$, then $p(X)$ has a well-defined finite degree $m := \deg(p) \ge 0$. But every non-zero polynomial in $U_{m+1}$ has degree at least $m + 1$, which implies $p(X) \notin U_{m+1}$, a contradiction. Hence, no non-zero polynomial belongs to the intersection, proving that
    \[
        \bigcap_{n=0}^\infty U_n = \{0\}.
    \]
\end{enumerate}
\exinfo{Official Solution (Problem Sheet 9, Exercise 6). Proves stabilization of descending subspace chains in finite dimensions and provides polynomial counterexamples in infinite dimensions.}
\end{exercisesolution}


